\section{背景介绍}

\label{sec:introduction}

\subsection{问题背景与数据来源}

乳腺癌的精确诊断与分类在临床医学中至关重要。本项目的核心分析数据来自于临床蛋白组学肿瘤分析财团 (NCI/NIH) 发布的 77 例乳腺癌样本的 iTRAQ 蛋白质组学数据。该数据集包含了每个样本中超过 12,000 种蛋白质的表达值。

\subsection{生物学意义}

在细胞运作机制中，DNA 承载的基因信息被转录为 RNA，进而翻译成蛋白质。由于蛋白质是负责细胞分裂、DNA 修复以及信号传导等功能的主要执行单元，DNA 突变最终会深刻改变癌细胞内的蛋白质表达图谱。因此，这些高维的蛋白质表达值与乳腺癌的特性、肿瘤几何形态以及临床分类结果息息相关。

\subsection{研究现状与项目目标}

在原始研究中，科学家通过对这些蛋白质数据进行 K-means 聚类分析，成功将乳腺癌患者划分为具有独特蛋白质表达特征的 3 个亚型。基于这一前沿背景，本课程项目旨在完整走通数据科学的分析流程：首先对存在缺失值的蛋白质高维特征进行预处理；随后，分别通过无监督的聚类分析和有监督的分类算法对数据进行深度挖掘；最终目标是评估这些蛋白质表达特征与现有临床标准（如 PAM50 mRNA）之间的关联，并尝试筛选出具有关键分类价值的特征属性。
